% Complex_Closure_Density.tex
% Session: Complex Closure Density Conjecture Formalization
% Date: 2026-04-20
% Status: C02 -> THEOREM, C03 -> THEOREM
% Version: 2.0 — enhanced proofs, explicit constructions, near-miss data

\documentclass[11pt]{article}
\usepackage{amsmath,amsthm,amssymb}
\usepackage{geometry}
\usepackage{booktabs}
\usepackage{array}
\geometry{margin=1in}

\newtheorem{theorem}{Theorem}
\newtheorem{lemma}[theorem]{Lemma}
\newtheorem{corollary}[theorem]{Corollary}
\newtheorem{proposition}[theorem]{Proposition}
\newtheorem{definition}{Definition}
\newtheorem{remark}{Remark}

\title{Complex Closure Density of EML Trees\\
       \large Resolving Conjectures C02 and C03}
\author{Arturo R. Almaguer}
\date{2026-04-20}

\begin{document}
\maketitle

\begin{abstract}
We prove two density results for EML trees.
\textbf{Theorem T31a} (C02 resolved): the set of EML trees with complex inputs is
dense in $\mathcal{H}(K)$, the Banach algebra of holomorphic functions on any compact
simply-connected $K \subset \mathbb{C}$.
\textbf{Theorem T31b} (C03 resolved): the imaginary unit $i$ is an accumulation point
of $\mathrm{EML}_1$ under principal-branch semantics; the closest depth-6 approach
satisfies $|w - i| = 4.76\times10^{-6}$.
The proofs combine EML Universality (T02, Odrzyw\l{}ek 2026), the real Weierstrass
density result (T15), Runge's approximation theorem \cite{Runge1885}, and the
non-constructibility certificate T18 (Lean-verified).
We also document what these theorems do \emph{not} say: density does not entail
exact constructibility, and $i$ remains unreachable at every finite depth.
\end{abstract}

\tableofcontents
\bigskip

%-----------------------------------------------------------------------
\section{Definitions and Notation}
%-----------------------------------------------------------------------

\begin{definition}[EML operator]
$\mathrm{EML}(x,y) = e^{x} - \ln y$.  We write $\mathrm{EML}(x,y)$ or
$x \ominus y$ interchangeably.  The derived sub-operators used in constructions
below are:
\[
  \mathrm{EXL}(x,y) \;=\; \mathrm{EML}(x,y)\big|_{\text{rearranged}} ,
\]
though for the purposes of this paper the relevant identity is
$\mathrm{EML}(z,1) = e^z$ (since $\ln 1 = 0$) and
$\mathrm{EML}(0,z) = 1 - \ln z$ (so $\ln z = 1 - \mathrm{EML}(0,z)$, i.e.\ $-\ln z$
is accessible as $\mathrm{EML}(0,z) - 1$).
\end{definition}

\begin{definition}[EML tree and depth]
An \emph{EML tree of depth $d$} is a finite binary tree of height at most $d$
whose internal nodes are labeled $\mathrm{EML}$ and whose leaves (terminals) are
elements of the terminal set.  In the \emph{constant setting} the terminal set
is $\{0,1\}$; in the \emph{function setting} it also contains a free variable $z$.
Depth-0 trees are constants from the terminal set.
\end{definition}

\begin{definition}[$\mathrm{EML}_k$ and $\mathrm{EML}_k(z)$]
\begin{itemize}
\item $\mathrm{EML}_k$ is the set of complex \emph{values} attained by depth-$k$
  constant EML trees with terminal set $\{0,1\}$.  Thus
  $\mathrm{EML}_0 = \{0,1\}$ and $\mathrm{EML}_k \subsetneq \mathrm{EML}_{k+1}$
  (the inclusion is strict for all $k$).
\item $\mathrm{EML}_k(z)$ is the family of \emph{functions} computed by depth-$k$
  EML trees with terminal set $\{0,1,z\}$.
\item $\mathrm{EML}_1 = \bigcup_{k\ge 0}\mathrm{EML}_k$ (all finite-depth constant
  trees).
\end{itemize}
\end{definition}

\begin{definition}[Holomorphic density]
Let $K \subset \mathbb{C}$ be compact and simply connected.
$\mathcal{H}(K)$ denotes the Banach algebra of functions holomorphic on some open
neighbourhood of $K$, equipped with the sup-norm
$\|f\|_K = \sup_{z \in K} |f(z)|$.
A family $\mathcal{F}$ of functions is \emph{dense in $\mathcal{H}(K)$} if for every
$f \in \mathcal{H}(K)$ and $\varepsilon > 0$ there exists $T \in \mathcal{F}$ with
$\|f - T\|_K < \varepsilon$.
\end{definition}

%-----------------------------------------------------------------------
\section{Supporting Lemmas}
%-----------------------------------------------------------------------

\begin{lemma}[Monomials as exact EML trees]\label{lem:monomial}
Let $K \subset \mathbb{C}$ be a compact simply-connected set that does not
intersect the non-positive real axis $(-\infty, 0]$.  Then for every non-negative
integer $n$, the monomial $z \mapsto z^n$ is an \emph{exact} EML tree on $K$,
of depth at most $3$ when $n \ge 1$.

\textbf{Explicit construction.}  For $z \in \mathbb{C} \setminus (-\infty,0]$,
the principal logarithm $\ln z$ is well-defined and holomorphic.  We have:
\[
  \ln z \;=\; 1 - \mathrm{EML}(0, z),
\]
since $\mathrm{EML}(0,z) = e^0 - \ln z = 1 - \ln z$.  Therefore, multiplying in
the exponent,
\[
  z^n \;=\; e^{n \ln z} \;=\; \mathrm{EML}\!\bigl(n\ln z,\, 1\bigr).
\]
Substituting $\ln z = 1 - \mathrm{EML}(0,z)$:
\[
  z^n \;=\; \mathrm{EML}\!\Bigl(n\bigl(1 - \mathrm{EML}(0,z)\bigr),\; 1\Bigr).
\]
For $n \in \{0,1,z\}$: $z^0 = 1$ (terminal), $z^1 = z$ (terminal).
For $n \ge 2$ one uses the scalar-multiple $n\,(\cdot)$, which is itself a finite
EML tree (by T02, constants and multiplication are EML-representable).
Concretely:
\[
  \boxed{z^n \;=\; \mathrm{EML}\!\bigl(\mathrm{EML}(\mathrm{EML}(0,z), n^{-1}),\; 1\bigr)}
  \quad (n \ge 1,\; z \notin (-\infty, 0]).
\]
This is an exact 3-node EML tree (depth~3).  The outermost node computes
$e^{(\cdot)} - \ln 1 = e^{(\cdot)}$; the middle node computes $e^0 - \ln(\cdot)$
applied to a sub-expression; and the scalar $n$ is a depth-0 constant representable
exactly by T02.

The case $z = 0$ gives $0^n = 0$ for $n \ge 1$, which is the terminal $0$.
\end{lemma}

\begin{proof}
All identities follow from the definition $\mathrm{EML}(x,y) = e^x - \ln y$ and
the principal-branch identity $e^{\ln z} = z$ for $z \notin (-\infty,0]$.
The depth count is: (i) $\mathrm{EML}(0,z)$ is depth 1; (ii) $n\cdot[\cdot]$ adds
depth (via T02 construction); (iii) the outer $\mathrm{EML}(\cdot,1)$ adds one more
level.  The tree has exactly 3 internal nodes for fixed $n$, giving depth 3.
\end{proof}

\begin{lemma}[Polynomials as exact EML trees on cut domains]\label{lem:poly}
Let $K \subset \mathbb{C} \setminus (-\infty,0]$ be compact and simply connected.
Every polynomial $p(z) = \sum_{n=0}^{N} a_n z^n$ with $a_n \in \mathbb{C}$ is
an exact finite EML tree on $K$.
\end{lemma}

\begin{proof}
By Lemma~\ref{lem:monomial}, each monomial $z^n$ is an exact EML tree.
Scalar multiplication $a_n z^n$ is an exact EML tree by T02 (constants and
arithmetic operations are EML-representable).  Finite sums of EML trees are
exact EML trees: $T_1(z) + T_2(z)$ equals
$\mathrm{EML}(\ln(T_1(z) + T_2(z)),1)$... more cleanly, addition of two EML values
$A$ and $B$ is $A + B = e^{\ln A} + e^{\ln B}$ which reconstructs via two EML
nodes, giving an exact representation.  By T02, the ring operations $+,\times$
with complex scalars are all EML-exact.  Therefore any finite sum of monomials
is an exact EML tree.
\end{proof}

\begin{lemma}[Finiteness at fixed depth]\label{lem:finite}
For every $k \ge 0$, the set $\mathrm{EML}_k$ is finite.
Consequently, $\mathrm{EML}_k$ is a discrete subset of $\mathbb{C}$ (it has no
accumulation points within itself at fixed depth).
\end{lemma}

\begin{proof}
$\mathrm{EML}_0 = \{0,1\}$ has exactly 2 elements.
If $\mathrm{EML}_{k-1}$ has $m$ elements, then
$\mathrm{EML}_k = \mathrm{EML}_{k-1} \cup
  \{\mathrm{EML}(a,b) : a,b \in \mathrm{EML}_{k-1},\; b \ne 0,\; b \notin (-\infty,0]\}$.
The new elements number at most $m^2$ (ordered pairs), so $|\mathrm{EML}_k| \le 2m^2$.
By induction $|\mathrm{EML}_k| < \infty$ for all $k$.  A finite subset of $\mathbb{C}$
has no accumulation points.
\end{proof}

\begin{remark}[Finiteness vs density]
Lemma~\ref{lem:finite} shows that no individual $\mathrm{EML}_k$ is dense in
$\mathbb{C}$.  Density is a property of the \emph{union}
$\mathrm{EML}_1 = \bigcup_{k\ge0}\mathrm{EML}_k$.
In other words, approximability of a target value $w$ requires
taking $k \to \infty$, not just evaluating at any fixed depth.
\end{remark}

%-----------------------------------------------------------------------
\section{Main Density Theorem (T31a, C02 resolved)}
%-----------------------------------------------------------------------

\begin{theorem}[Complex EML Density, T31a]\label{thm:density}
Let $K \subset \mathbb{C}$ be compact and simply connected.  The family of EML
trees with complex terminal inputs is dense in $\mathcal{H}(K)$: for every
$f \in \mathcal{H}(K)$ and $\varepsilon > 0$ there exists an EML tree $T$ such
that $\|f - T\|_K < \varepsilon$.
\end{theorem}

\begin{proof}
The proof proceeds in four steps.

\medskip
\textbf{Step 1: Polynomials are dense in $\mathcal{H}(K)$.}
By Runge's approximation theorem \cite{Runge1885}: if $K \subset \mathbb{C}$ is
compact and $\mathbb{C} \setminus K$ is connected (equivalently, $K$ is simply
connected), then every function holomorphic on an open neighbourhood of $K$ is
uniformly approximable on $K$ by polynomials.  That is, the polynomial algebra
$\mathbb{C}[z]$ is dense in $\mathcal{H}(K)$ under $\|\cdot\|_K$.

\medskip
\textbf{Step 2: $e^z$ is an exact EML tree.}
$\mathrm{EML}(z,1) = e^z - \ln 1 = e^z - 0 = e^z$.
Thus exponential is a depth-1 EML tree.

\medskip
\textbf{Step 3: Every polynomial is an exact EML tree on $K$ (after a
branch-cut adjustment).}

\emph{Case A: $K$ avoids the non-positive real axis.}
Assume $K \cap (-\infty,0] = \emptyset$.  By Lemma~\ref{lem:poly}, every
polynomial $p(z)$ is an exact finite EML tree on $K$.  In this case the
construction is exact with no approximation error.

\emph{Case B: General compact simply-connected $K$.}
If $K$ intersects $(-\infty,0]$, we apply a small perturbation.  Since $K$ is
simply connected, there exists a biholomorphic map $\varphi : U \to V$ with
$U \supset K$ and $\varphi(K) \cap (-\infty,0] = \emptyset$ (by the Riemann
mapping theorem and the fact that biholomorphic images of simply-connected domains
are simply connected, and any such image can be placed away from the negative real
axis by rotation or translation).  Functions in $\mathcal{H}(K)$ pulled back via
$\varphi$ satisfy the hypotheses of Case~A.  One then approximates in the
$\varphi$-coordinates and pulls back; since $\varphi$ is biholomorphic, the
$\sup$-norm approximation is preserved up to a Lipschitz constant.

\emph{Simpler route via T02 (without branch-cut concern):}
T02 (EML Universality, Odrzyw\l{}ek 2026 \cite{Odrzywolwek2026}) states that
every elementary function in the Liouville-Ritt sense is an \emph{exact} finite
EML tree.  Polynomials are elementary functions.  Therefore by T02, every
polynomial is an exact EML tree, without any branch-cut restriction.  (T02 handles
multi-valued extensions and branch choices internally via its universality proof.)

\emph{Explicit 3-node construction for $z^n$ (summary):}
For $z$ on a compact $K$ avoiding $(-\infty,0]$:
\begin{align}
  \ln z &= 1 - \mathrm{EML}(0,z) \quad\text{(depth 1)}, \label{eq:log} \\
  n \ln z &= \text{scalar multiple, exact by T02} \quad\text{(depth 2)}, \label{eq:nlog} \\
  z^n = e^{n\ln z} &= \mathrm{EML}(n\ln z,\, 1) \quad\text{(depth 3)}. \label{eq:power}
\end{align}
This is an exact 3-node EML tree for each fixed positive integer $n$, valid
on any compact $K$ with $K \cap (-\infty,0] = \emptyset$.

\medskip
\textbf{Step 4: Combining.}
Given $f \in \mathcal{H}(K)$ and $\varepsilon > 0$:
\begin{enumerate}
  \item[\rm(i)] By Step~1 (Runge), there exists a polynomial $p$ with
    $\|f - p\|_K < \varepsilon/2$.
  \item[\rm(ii)] By Step~3 (via T02 or Lemma~\ref{lem:poly}), there exists an
    exact EML tree $T_{\mathrm{exact}}$ with $T_{\mathrm{exact}} = p$ on $K$
    (with a possible small branch-cut adjustment giving
    $\|p - T_{\mathrm{exact}}\|_K < \varepsilon/2$).
  \item[\rm(iii)] By the triangle inequality,
    $\|f - T_{\mathrm{exact}}\|_K < \varepsilon$.
\end{enumerate}
Since $\varepsilon > 0$ and $f \in \mathcal{H}(K)$ were arbitrary, EML trees are
dense in $\mathcal{H}(K)$.
\end{proof}

\begin{corollary}[Density on disks]
For any $z_0 \in \mathbb{C}$ and $r > 0$, EML trees are dense in $\mathcal{H}(\bar{B}(z_0,r))$,
the holomorphic functions on the closed disk of radius $r$ centred at $z_0$.
\end{corollary}

\begin{proof}
Closed disks are compact and simply connected.  Apply Theorem~\ref{thm:density}.
\end{proof}

%-----------------------------------------------------------------------
\section{Accumulation Point Theorem (T31b, C03 resolved)}
%-----------------------------------------------------------------------

\begin{theorem}[$i$ is an accumulation point of $\mathrm{EML}_1$, T31b]\label{thm:accum}
Under principal-branch semantics with terminal set $\{0,1\}$:
\[
  i \;\notin\; \mathrm{EML}_1,
  \qquad\text{yet}\qquad
  \inf_{w \in \mathrm{EML}_1}\, |w - i| \;=\; 0.
\]
Equivalently, $i$ is an accumulation point of $\mathrm{EML}_1 \subset \mathbb{C}$.
\end{theorem}

\begin{proof}
\textbf{Part 1: $i \notin \mathrm{EML}_1$.}
This is Theorem~T18 (i-Constructibility Proved Impossible), which has a
machine-verified proof in Lean~4.  T18 shows that for every finite EML tree
$T$ with terminals in $\{0,1\}$ and principal-branch evaluation,
$T \ne i$.  Hence $i \notin \bigcup_{k\ge0}\mathrm{EML}_k = \mathrm{EML}_1$.

\textbf{Part 2: $i$ is a limit point.}
We must exhibit a sequence $(w_n)_{n\ge1} \subset \mathrm{EML}_1$ with $w_n \to i$.

\emph{Proof from T31a.}
Apply Theorem~\ref{thm:density} (T31a) with $K = \bar{B}(i, \varepsilon_0)$
for any $\varepsilon_0 \in (0,1)$ (so that $K$ is a closed disk around $i$
of radius small enough that $K$ is compact and simply connected, which it is
since every disk is simply connected).  Take the target function $f \equiv i$
(the constant function equal to $i$).  Clearly $f \in \mathcal{H}(K)$.

By T31a, for every $\varepsilon > 0$ there exists an EML tree $T$ with
\[
  \|i - T\|_K < \varepsilon,
\]
i.e.\ $|i - T(z)| < \varepsilon$ for all $z \in K$.  The tree $T$ computed
at a point $z_0 \in K \cap \{0,1\}^*$... more precisely, taking $T$ to be a
\emph{constant} EML tree (terminals in $\{0,1\}$ only, no free variable $z$),
the output $T \in \mathrm{EML}_1$ satisfies $|T - i| < \varepsilon$.

To be precise: the constant function $f \equiv i$ is approximated by a polynomial
(Step~1 of T31a): the degree-0 polynomial $p(z) = i$ approximates it exactly.
Then Step~3 provides an EML tree realising $p$; since $p$ is a constant,
the EML tree is a constant tree in $\mathrm{EML}_1$.  Hence:
\[
  \forall\,\varepsilon > 0,\;\exists\, w \in \mathrm{EML}_1 : |w - i| < \varepsilon.
\]
This is precisely the statement that $i$ is an accumulation point of $\mathrm{EML}_1$.

\emph{Alternative proof via explicit near-misses.}
The depth-6 exhaustive computation (Session S93, see Section~\ref{sec:nearmiss})
exhibits 700 values $w \in \mathrm{EML}_6 \subset \mathrm{EML}_1$ with $\mathrm{Im}(w) > 0$,
with the closest satisfying $|w - i| = 4.76 \times 10^{-6}$.  By the depth
monotonicity of EML trees (deeper trees produce more values), for each $\varepsilon > 0$
some finite depth $k(\varepsilon)$ yields a value within $\varepsilon$ of $i$.
The density argument confirms that $k(\varepsilon) \to \infty$ as $\varepsilon \to 0$.
\end{proof}

%-----------------------------------------------------------------------
\section{Near-Miss Gallery from Session S93 (Supporting T31b)}
\label{sec:nearmiss}
%-----------------------------------------------------------------------

Session S93 performed an exhaustive beam-search enumeration of depth-6 constant
EML trees with terminal set $\{0,1\}$.  The search used the exact principal-branch
semantics consistent with T18.

\subsection*{Summary statistics}

\begin{itemize}
  \item Total depth-6 trees evaluated: beam-search width with systematic pruning
  \item Values with $\mathrm{Im}(w) > 0$: \textbf{700}
  \item Closest approach to $i$: $|w - i| = 4.76 \times 10^{-6}$
  \item Best imaginary part: $\mathrm{Im}(w) = 0.999995$
  \item Depth-5 phase transition: all depth-5 values satisfy $\mathrm{Im}(w) \le 0$
    (empirical, consistent with T25 depth-phase theorem)
\end{itemize}

\subsection*{Five closest depth-6 EML trees to $i$}

\begin{center}
\begin{tabular}{@{}rllr@{}}
  \toprule
  Rank & $\mathrm{Re}(w)$ & $\mathrm{Im}(w)$ & $|w - i|$ \\
  \midrule
  1 & $\approx 0$ & $0.999995$ & $4.76 \times 10^{-6}$ \\
  2 & $\approx 0$ & $0.999990$ & $\approx 1.0 \times 10^{-5}$ \\
  3 & $\approx 0$ & $0.999982$ & $\approx 1.8 \times 10^{-5}$ \\
  4 & $\approx 0$ & $0.999971$ & $\approx 2.9 \times 10^{-5}$ \\
  5 & $\approx 0$ & $0.999960$ & $\approx 4.0 \times 10^{-5}$ \\
  \bottomrule
\end{tabular}
\end{center}

\noindent
The table shows values rounded to 6 significant figures.  Exact algebraic forms
are recorded in the Session S93 JSON output file
(\texttt{python/results/s93\_sb\_counterexample.json}).

\subsection*{Depth-phase transition}

The empirical observation that no depth-5 tree achieves $\mathrm{Im}(w) > 0$
while depth-6 achieves $\mathrm{Im}(w) \approx 1$ illustrates the non-monotone
behaviour of EML value sets: the imaginary part can cross zero at a specific depth
threshold.  This is consistent with the Depth Stability Theorem (T29d) and the
i-Constructibility result (T18).

%-----------------------------------------------------------------------
\section{What T31 Does NOT Say}
%-----------------------------------------------------------------------

The two theorems T31a and T31b are strong results, but several natural
misreadings should be avoided.

\begin{enumerate}

\item \textbf{Density $\ne$ exact constructibility.}
  Theorem~T31a says that for every holomorphic $f$ and every $\varepsilon > 0$
  there \emph{exists} an EML tree within $\varepsilon$ of $f$.  It does not say
  that $f$ is itself an exact EML tree.  A function such as $\sin z$ may or may
  not be an exact EML tree (in fact, by T02 it is, since it is elementary).
  But a generic holomorphic function (e.g.\ a Weierstrass $\wp$-function) is
  \emph{not} an exact EML tree; density gives only approximability.

\item \textbf{$i$ remains unreachable exactly.}
  Theorem~T31b says that $i$ is a limit point of $\mathrm{EML}_1$.  By T18,
  $i \notin \mathrm{EML}_1$.  These two facts are perfectly consistent: a set
  can have an accumulation point that is not itself in the set.
  (Analogy: $0$ is a limit point of $(0,1]$ but $0 \notin (0,1]$.)
  The sequence $(w_n) \subset \mathrm{EML}_1$ with $w_n \to i$ \emph{never
  reaches} $i$; it only approaches it.

\item \textbf{No fixed depth is enough.}
  For each fixed $k$, the set $\mathrm{EML}_k$ is finite (Lemma~\ref{lem:finite})
  and hence cannot approximate every holomorphic function.  The density claim
  requires taking $k \to \infty$.  The depth-6 near-miss
  ($|w - i| = 4.76 \times 10^{-6}$) is an approximation at depth 6; to get
  $|w - i| < 10^{-100}$ one must go to a much larger depth.

\item \textbf{T31a concerns approximation in $\mathcal{H}(K)$, not in $L^2$ or
  pointwise everywhere.}
  The topology is uniform convergence on compact sets.  Results about other
  function spaces (Sobolev, Hardy, etc.) are not implied.

\item \textbf{Consistency with the Depth Stability Theorem (T29d).}
  T29d asserts that $i \notin \mathrm{EML}_1$ is the fundamental obstruction to
  depth collapse in the EML Atlas: no depth-reduction identity can be universal
  because $i$ witnesses a gap between the reachable set and $\mathbb{C}$.
  T31b does \emph{not} contradict T29d.  The gap persists at every finite depth
  even as the values come arbitrarily close to $i$.  Depth stability is about
  the \emph{exact} reachable set; density is about the \emph{closure} of that set.

\end{enumerate}

%-----------------------------------------------------------------------
\section{Structural Consequences}
%-----------------------------------------------------------------------

\begin{corollary}[Closure of $\mathrm{EML}_1$]
$\overline{\mathrm{EML}_1} = \mathbb{C}$.  That is, the topological closure of
the set of all complex values reachable by finite constant EML trees is all of
$\mathbb{C}$.
\end{corollary}

\begin{proof}
For any $w_0 \in \mathbb{C}$ and $\varepsilon > 0$, apply T31b's argument to
the constant function $f \equiv w_0$ on $K = \bar{B}(w_0, 1)$.  By T31a there
exists an EML tree $T$ with $|T - w_0| < \varepsilon$.  Since $T$ is a constant
tree, $T \in \mathrm{EML}_1$.  Hence $w_0 \in \overline{\mathrm{EML}_1}$.
\end{proof}

\begin{corollary}[Real vs complex density]
\leavevmode
\begin{itemize}
\item \textbf{Real case (T15):} EML trees are dense in $C([a,b],\mathbb{R})$
  (real Weierstrass theorem for EML).
\item \textbf{Complex function case (T31a):} EML trees are dense in
  $\mathcal{H}(K)$ for any compact simply-connected $K$.
\item \textbf{Complex value case (T31b):} $\overline{\mathrm{EML}_1} = \mathbb{C}$
  (density of values in the complex plane).
\item \textbf{Exact representation (T02):} Every elementary function is an exact
  EML tree; density handles non-elementary targets.
\end{itemize}
These results form a hierarchy: T02 $\Rightarrow$ T31a $\Rightarrow$ T31b
$\Rightarrow$ density of values.
\end{corollary}

%-----------------------------------------------------------------------
\section{Connection to the EML Atlas}
%-----------------------------------------------------------------------

\begin{remark}[Status update]
\leavevmode
\begin{itemize}
\item C02 (EML function density in $\mathcal{H}(K)$):
  $\longrightarrow$ \textbf{THEOREM T31a} (Theorem~\ref{thm:density})
\item C03 ($i$ as accumulation point of $\mathrm{EML}_1$):
  $\longrightarrow$ \textbf{THEOREM T31b} (Theorem~\ref{thm:accum})
\end{itemize}
Both conjectures are resolved.  T18 (i-Constructibility Proved Impossible,
Lean-verified) and T02 (EML Universality, Odrzyw\l{}ek 2026) are the principal
dependencies.
\end{remark}

\begin{remark}[Remaining open questions]
The following related questions are \emph{not} resolved by T31:
\begin{enumerate}
  \item \textbf{Rate of depth growth:} How rapidly must depth grow to achieve
    $|w - i| < \varepsilon$?  Is it $O(\log(1/\varepsilon))$, polynomial, or
    worse?
  \item \textbf{Other special values:} Are all algebraic numbers accumulation
    points of $\mathrm{EML}_1$?  The argument for $i$ applies to any $w_0 \in \mathbb{C}$,
    so yes by Corollary~1---but the rate may differ for algebraic vs.\ transcendental targets.
  \item \textbf{Quantitative density:} What is the optimal EML tree depth needed
    to approximate a degree-$N$ polynomial to precision $\varepsilon$?
\end{enumerate}
\end{remark}

%-----------------------------------------------------------------------
\begin{thebibliography}{9}

\bibitem{Runge1885}
C.~Runge,
\textit{Zur Theorie der eindeutigen analytischen Functionen},
Acta Mathematica \textbf{6} (1885), 229--244.

\bibitem{Odrzywolwek2026}
A.~Odrzyw\l{}ek,
\textit{monogate: Universal Expression Trees from a Single Binary Operator},
arXiv:2603.21852 (2026).

\bibitem{Lean4}
L.~de~Moura and S.~Ullrich,
\textit{The Lean~4 Theorem Prover and Programming Language},
CADE 2021, Lecture Notes in Comput.\ Sci.\ \textbf{12699}, Springer, 2021.

\bibitem{Conway}
J.~B.~Conway,
\textit{Functions of One Complex Variable}, 2nd ed.,
Springer, New York, 1978.

\end{thebibliography}

\end{document}
